\documentclass[11pt,letterpaper]{article}
\usepackage[T1]{fontenc}
\usepackage[utf8]{inputenc}
\usepackage{lmodern}
\usepackage[margin=1in]{geometry}
\usepackage{amsmath,amssymb,amsthm,mathtools}
\usepackage{microtype,booktabs,array,longtable,xcolor,listings,fancyhdr,xurl}
\usepackage[hidelinks]{hyperref}
\definecolor{ink}{HTML}{17324D}
\hypersetup{pdftitle={Smaller clause-discharge checkers for iPRA},
 pdfauthor={Research prepared for Vladimir Reshetnikov},
 pdfsubject={A complete smaller implementation, deterministic witness discharge, and exact bi-interpretability}}
\newtheorem{theorem}{Theorem}[section]
\newtheorem{lemma}[theorem]{Lemma}
\newtheorem{proposition}[theorem]{Proposition}
\newtheorem{corollary}[theorem]{Corollary}
\theoremstyle{definition}\newtheorem{definition}[theorem]{Definition}
\theoremstyle{remark}\newtheorem{remark}[theorem]{Remark}
\newcommand{\nf}{\operatorname{nf}}
\newcommand{\ipra}{\mathsf{iPRA}}
\newcommand{\FV}{\operatorname{FV}}
\newcommand{\D}{\mathsf D}
\newcommand{\Rcl}{\mathcal R}
\newcommand{\one}{\overline{1}}
\newcommand{\two}{\overline{2}}
\setlength{\parskip}{3pt}
\setlength{\emergencystretch}{3em}
\setlength{\headheight}{14pt}
\pagestyle{fancy}\fancyhf{}
\fancyhead[L]{\small Smaller clause-discharge checkers for iPRA}
\fancyhead[R]{\small November 11, 2026}
\fancyfoot[C]{\thepage}
\lstset{language=Python,basicstyle=\ttfamily\scriptsize,
 columns=fullflexible,keepspaces=true,breaklines=true,
 showstringspaces=false,frame=single,rulecolor=\color{black!20},
 xleftmargin=4pt,xrightmargin=4pt,aboveskip=9pt,belowskip=9pt}
\title{\color{ink}Smaller clause-discharge checkers for iPRA\\[6pt]
\large A measured implementation reduction and a deterministic equality rule}
\author{Research prepared for Vladimir Reshetnikov}
\date{November 11, 2026}
\begin{document}
\maketitle
\begin{abstract}
We give two complete presentations of full first-order intuitionistic
primitive recursive arithmetic over the same normal scalar iterator syntax.
The implementation recommended when checker size is decisive retains the
existing three-rule witness-discharge calculus and reduces its complete
normalized Python checker from 2,273 to 2,010 bytes. The second presentation
has a new deterministic discharge rule: the checker reads an equality from
the opened witness clause at the fixed argument $1$, so the certificate
supplies neither an equality proof nor a substitution target. It uses the
same three inference families and the same single reflexivity axiom clause.
Its complete checker has 2,058 bytes, 653 significant tokens and 729 AST
nodes. The smaller explicit checker has 636 tokens and 708 AST nodes.

We prove exact equivalence on arbitrary well-scoped clause sequents. The
noncircular completeness argument first derives numerical cases without
equality transport, and then uses an inherited two-row selector package
with the equality row moved to the fixed inspection position. Both
presentations inherit whole-domain bi-interpretability and equiconsistency
with iPRA, actual equality, and identity numerical composite maps.

The deterministic rule costs 48 bytes more than the matched explicit rule;
its benefit is a smaller primitive certificate interface. The implementation
reductions and the logical reduction are therefore reported separately.
Neither checker beats the repository's distinct 516-byte compact codec or
788-byte validated public-goal codec, and no absolute minimum is claimed.
The package contains the complete source, mathematical foundations,
executable proof translations, adversarial checks, and a bounded-certificate
variant. New proofs are written mathematics, not a new Lean formalization.
\end{abstract}
\clearpage
\begingroup\fontsize{9}{10}\selectfont\setlength{\parskip}{0pt}\setcounter{tocdepth}{1}
\tableofcontents\endgroup
\clearpage

\section{Result and the comparison that the measurements support}
\label{sec:result}

The request is to simplify the current proposal at
\path{Logic/PrimitiveRecursiveArithmetic/ConciseIPRA}. We inspected the
repository at the fixed revision
\texttt{ffa94f76bc01}, identified in full in the bibliography, including the six
archives in \path{docs/incoming/}. That snapshot matters: the incoming
witness-discharge proposal already has three inference families, although
the main README's most compact clause interface has four families.
The closest logical predecessor is thus the incoming witness-discharge
package~\cite{prior}, not an earlier eight-rule native presentation.

Two distinct simplifications result. First, a complete implementation can
infer the next fresh numerical parameter from its clause context, use one
shared traversal for shifting and opening, and expose a minimal two-field
standalone input. This preserves the predecessor's calculus and lowers all
three primary source coordinates. Second, the equality proof in witness
discharge can be replaced by a fixed, syntactically recognizable consequence
of the opened clause. This lowers the rule's proof-premise count from three
to two and removes a term parameter, with no loss of arithmetic or logic.
Its selector must be counted, however, and it makes the matched program
slightly longer.

\subsection{Which delivered presentation should be used?}

If the objective is the size of a straightforward complete proof checker,
use \path{code/checker_explicit.py}. It is the smallest measured member of
the delivered pair. If the objective is to minimize the equality information
that a proof producer supplies, use
\path{code/checker_deterministic.py}. The latter is the new mathematical
calculus. Both have the same numerical and formula language, the same
ordinary meaning, and the same consequences under every canonical clause
context. The explicit implementation is included as a full deliverable,
not merely as a control used to make a comparison favorable.

\begin{center}
\small
\begin{tabular}{@{}lrrrr@{}}
\toprule
Complete standalone module & Bytes & Tokens & AST & Gzip-9\\
\midrule
Prior witness discharge & 2273 & 732 & 813 & 903\\
Prior uniform binding & 2238 & 693 & 764 & 866\\
Delivered explicit discharge & 2010 & 636 & 708 & 791\\
Delivered deterministic discharge & 2058 & 653 & 729 & 805\\
\bottomrule
\end{tabular}
\end{center}

These counts include the parser invocation, typing, scopes, normalizer,
capture-avoiding substitution, every inference, the complete input protocol,
and final acceptance. They are reproduced by the supplied measurement
script on CPython 3.12.14: remove documentation strings, apply
\texttt{ast.unparse}, append one LF, count UTF-8 bytes, significant lexical
tokens, and all nodes visited by \texttt{ast.walk}. Gzip uses level 9 and
timestamp zero and is informational, not an objective being minimized.
The exact raw-source sizes and hashes are also recorded.

The explicit checker saves 263 normalized bytes (11.57\%), 96 tokens
(13.11\%), and 105 AST nodes (12.92\%) against its direct predecessor.
The deterministic checker saves 215 bytes (9.46\%), 79 tokens (10.79\%),
and 84 nodes (10.33\%) against that predecessor. In the matched delivered
pair, automatic extraction adds 48 bytes, 17 tokens, 21 nodes and 14 gzip
bytes. Thus the theorem about the new rule is not an explanation of all
the implementation savings.

\subsection{Scope of the minimum claim}

The inspected readable/transparent incoming checkers have coordinate
minima 2,238 bytes, 666 tokens and 764 AST nodes, achieved by different
programs~\cite{binding,clauses}. Both delivered programs are below those
three numbers. This is a scoped comparison, not a lower-bound theorem.
The repository also contains heavily fused compact programs, including
a 516-byte checker whose public-goal grammar is a precondition and a
788-byte native public-goal validator. Those are smaller than this delivery.
There is no justified claim here to have found the globally simplest
axiomatization, or even the shortest checker across all codecs in the
repository. The positive result is a genuine logical restriction together
with smaller complete implementations in the explicitly compared family.

\subsection{What is inherited and what is new}

Inherited components are the scalar iterator arithmetic, the paired
$D$ binder, the witness-dependent introduction, optional-premise clauses,
the fixed reflexivity clause, full computation normalization, the two-row
selector construction, and the arithmetic bridge to iPRA. The new logical
step is to designate row $1$ in the elimination rule and read the equality
from it automatically, replacing the explicit equality subproof and target
term. The rows of the inherited selector package are interchanged to match
that fixed position. The new implementation work concerns the context-derived
scope, common traversal, and standalone protocol. All of these costs remain
in the measured modules.

The following sections state the complete calculus and prove the new step.
The arithmetic and logical foundations are reproduced with attribution in
the appendices. They are derived comparisons and do not enlarge the list of
primitive rules.
